\begin{tikzpicture}[
  font=\scriptsize,
  judul/.style={draw, rounded corners, fill=praditagreen!15, minimum height=0.55cm,
                minimum width=5.5cm, font=\scriptsize\bfseries},
  ket/.style={anchor=west, align=left, font=\scriptsize\bfseries},
  panah/.style={-{Stealth[length=2mm]}, thick},
  komposisi/.style={{Diamond[length=3mm, width=2mm]}-, thick},
  agregasi/.style={{Diamond[open, length=3mm, width=2mm]}-, thick},
  penugasan/.style={{Circle[length=1.6mm]}-{Stealth[length=2mm]}, thick},
  realisasi/.style={-{Triangle[open, length=2.2mm, width=2.4mm]}, thick, dashed},
  layani/.style={-{Straight Barb[length=2mm, width=2.4mm]}, thick},
  akses/.style={-{Straight Barb[length=1.4mm, width=1.8mm]}, thick, dotted},
  pengaruh/.style={-{Straight Barb[length=2mm, width=2.4mm]}, thick, dashed},
  asosiasi/.style={thick},
  pemicu/.style={-{Triangle[length=2.2mm, width=2.2mm]}, thick},
  aliran/.style={-{Triangle[length=2.2mm, width=2.2mm]}, thick, dashed},
  spesialisasi/.style={-{Triangle[open, length=2.5mm, width=2.8mm]}, thick}
]
  % Kolom kiri: struktural dan dinamis.
  \node[judul] at (2.75,0.0) {Struktural};
  \draw[komposisi] (0.0,-0.7) -- (1.5,-0.7);
  \node[ket] at (1.7,-0.7) {Komposisi\\{\tiny\mdseries bagian hanya milik satu induk}};
  \draw[agregasi] (0.0,-1.45) -- (1.5,-1.45);
  \node[ket] at (1.7,-1.45) {Agregasi\\{\tiny\mdseries bagian boleh dimiliki banyak induk}};
  \draw[penugasan] (0.0,-2.2) -- (1.5,-2.2);
  \node[ket] at (1.7,-2.2) {Penugasan\\{\tiny\mdseries siapa atau apa yang menjalankan}};
  \draw[realisasi] (0.0,-2.95) -- (1.5,-2.95);
  \node[ket] at (1.7,-2.95) {Realisasi\\{\tiny\mdseries yang konkret mewujudkan yang abstrak}};

  \node[judul] at (2.75,-3.75) {Dinamis};
  \draw[pemicu] (0.0,-4.45) -- (1.5,-4.45);
  \node[ket] at (1.7,-4.45) {Pemicu\\{\tiny\mdseries urutan sebab dan waktu}};
  \draw[aliran] (0.0,-5.2) -- (1.5,-5.2);
  \node[ket] at (1.7,-5.2) {Aliran\\{\tiny\mdseries sesuatu berpindah, tanpa urutan wajib}};

  % Kolom kanan: dependensi dan lainnya.
  \node[judul] at (9.05,0.0) {Dependensi};
  \draw[layani] (6.3,-0.7) -- (7.8,-0.7);
  \node[ket] at (8.0,-0.7) {Melayani\\{\tiny\mdseries penyedia menolong pemakai}};
  \draw[akses] (6.3,-1.45) -- (7.8,-1.45);
  \node[ket] at (8.0,-1.45) {Akses\\{\tiny\mdseries membaca atau menulis data}};
  \draw[pengaruh] (6.3,-2.2) -- (7.8,-2.2);
  \node[ket] at (8.0,-2.2) {Pengaruh\\{\tiny\mdseries mendorong atau menghambat tujuan}};
  \draw[asosiasi] (6.3,-2.95) -- (7.8,-2.95);
  \node[ket] at (8.0,-2.95) {Asosiasi\\{\tiny\mdseries terkait, tanpa arti khusus}};

  \node[judul] at (9.05,-3.75) {Lainnya};
  \draw[spesialisasi] (6.3,-4.45) -- (7.8,-4.45);
  \node[ket] at (8.0,-4.45) {Spesialisasi\\{\tiny\mdseries jenis khusus dari elemen lain}};
  \node[ket, font=\tiny\mdseries, gray] at (6.3,-5.2)
    {ujung bentuk belah ketupat dan bulatan ada di sisi induk\\atau pelaksana; ujung panah ada di sisi tujuan};
\end{tikzpicture}
